\section{术语与设计模式速查}

\subsection{核心术语对照}
为避免团队在跨模块协作时出现语义偏差，本讲可提炼出一套最小术语表。建议在项目文档中固定这些术语定义，减少 ``同名异义'' 导致的沟通成本与实验偏差。

\begin{table}[H]
\centering
\begin{tabular}{p{2.8cm}p{3.6cm}p{4.1cm}}
\toprule
\textbf{术语} & \textbf{课程语境中的含义} & \textbf{实现时的注意点} \\
\midrule
Reasoning action & 用于改进后续动作质量的中间思考步骤 & 需绑定触发条件和预算上限 \\
Reflection & 对既有轨迹进行自检并尝试修正 & 需要失败信号，否则易空转 \\
Replanning & 因环境偏差而重写计划 & 要明确重规划阈值，避免频繁抖动 \\
Working memory & 当前步骤所需的短时上下文 & 控制容量，避免挤占关键信息 \\
Episodic memory & 任务过程中的事件与结果记录 & 需附来源和时间戳，支持复盘 \\
Semantic memory & 稳定事实与知识库内容 & 管理版本与冲突，防止旧知识污染 \\
Procedural memory & 可复用策略、流程和模板 & 维护触发条件，避免错误泛化 \\
Verifier & 验证中间或终局正确性的机制 & 验证器本身也需持续评估 \\
\bottomrule
\end{tabular}
\caption{Reasoning、Memory、Planning 相关术语的工程化解释}
\end{table}

\subsection{常见设计模式}
课程实践可进一步沉淀为几种高频设计模式。它们不是银弹，但能显著降低从研究原型到生产系统的迁移摩擦。

\begin{knowledgebox}{四种常用设计模式}
\begin{enumerate}
    \item \textbf{Plan-Execute-Verify}：先计划、后执行、再验证，适合可验证任务；
    \item \textbf{Think-Act-Reflect}：推理与执行交替，适合动态环境；
    \item \textbf{Retrieve-Reason-Write}：先取证再推理，适合知识密集任务；
    \item \textbf{Shadow-to-Production}：先影子评测再灰度上线，适合高风险业务。
\end{enumerate}
\end{knowledgebox}

\begin{warningbox}{模式复用时的边界条件}
同一模式在不同任务分布下的效果可能相反。上线前应验证：任务可验证性、工具稳定性、延迟预算、权限约束是否满足模式前提。
\end{warningbox}

\subsection{本章小结}
统一术语和设计模式能减少团队协作摩擦，提高实验可比性与系统可维护性。这也是 Agent 工程从“个体经验”走向“组织能力”的关键一步。

